\subsection{Synthesis}
\label{section:fc_synthesis}

The chapter has traced the model-based reinforcement learning line from two 1990 origins, Sutton's Dyna and Schmidhuber's controller-model architecture, through the deep revival in \citet{HaSchmidhuber2018}, the recurrent state-space model that became the workhorse of the Dreamer line, the value-aware question that runs from VAML through value-equivalence to MuZero, the modern Dyna in continuous control associated with the MBPO ensemble line, and the convergence point at TD-MPC2. The remainder of the section notes where model-based methods win, where they fail, where the surrounding boundary makes the model unnecessary, and what remains open.

\subsubsection{Where model-based RL wins}
\label{section:fc_synthesis_wins}

Sample efficiency is the headline gain. Model-based methods outperform model-free baselines when environment interaction is scarce or expensive, including pixel-based continuous control in PlaNet and Dreamer, discrete-state planning in MuZero, and continuous control with sparse reward in MBPO and TD-MPC2. The learned model serves as a data amplifier through the Dyna planning step and as a representation amplifier, since the RSSM and TD-MPC2 latents transfer across tasks within their respective suites. When transferable representations or finite-data regimes are the binding constraint, model-based methods win.

\subsubsection{Where model-based RL fails}
\label{section:fc_synthesis_fails}

The first failure mode is misspecification and the objective mismatch between training and use. \citet{Lambert2020} document that maximum-likelihood model fitting and reward optimization are decoupled objectives, so a model trained to high likelihood is not necessarily a model whose induced policy is high-return. Their sweeps show likelihood-reward correlations as low as $0.07$ on Half-Cheetah and adversarial perturbations that preserve likelihood while halving return. \citet{VoelckerCalib2025} extend this point to MuZero and show that value-equivalent representations are systematically miscalibrated under distribution shift. The value-aware line (\S\ref{section:fc_value_aware}) is the partial answer, reformulating the training target around the value functional rather than the full distribution. A fully principled treatment of off-distribution calibration remains open.

The second failure mode is compounding error and the data-coverage gaps that exploration must fill. \citet{Asadi2018Lipschitz} give the local-operator Lipschitz bound for model-based RL, in which one-step prediction error compounds geometrically along the rollout at rate $K_F$, the Lipschitz constant of the learned dynamics. The induced value-error bound on the planner stays finite when $\gamma K_F < 1$ and diverges as $\gamma K_F \to 1$. When the contraction condition fails, model-based planning diverges. \citet{Talvitie2017} documents that retraining a single model on its own predictions can destabilize the model itself and proposes a self-correcting hallucinated-rollout scheme as one fix. \citet{Pathak2017} and \citet{Sekar2020} attack the same problem from the data side, with prediction-error-driven curiosity in the former and ensemble-disagreement-driven planning in the latter. Both push the agent's data distribution beyond what greedy exploitation alone produces, with measurable gains in regions of state space the agent would otherwise miss.

\subsubsection{When the model is not needed}
\label{section:fc_synthesis_boundary}

Not every sequential decision problem rewards the cost of learning a forecaster. When the decision reduces to a supervised prediction followed by a fixed downstream rule, a direct supervised reduction can match or beat a learned-world-model pipeline at much lower engineering cost, the boundary studied by \citet{Madeka2022DeepInventory} in deep inventory management. The boundary itself is closed-loop feedback. When the agent's action does not feed back into the data-generating process, a supervised reduction suffices. When it does, the learned model becomes load-bearing, because forecasts and decisions are then jointly determined. This is the boundary that places model-based reinforcement learning inside the sequential decision-making space and separates it from one-shot forecasting.

\subsubsection{Open questions}
\label{section:fc_synthesis_open}

The dual simulation in \S\ref{section:fc_dual_sim} maps where model-based reinforcement learning sits on the sample-efficiency frontier among the other paradigms in a single-agent cobweb with adjustment cost and a logistic-growth fishery. In the cobweb setting the Marcet-Sargent divergence-in-the-unstable-region story does not arise, because the environment is self-referential through prices alone and not through expectations. The more telling result is the order-of-magnitude separation between recursive least squares with correct functional form, a model-based learner that must estimate the demand and the cost coefficients, a population-based gradient-free search, and a tabular model-free agent. Whether the divergence story would reappear in an expectational cobweb in which beliefs about future prices drive supply, and in particular how a learned-neural-network world model would behave in that regime, is one of the open questions the chapter does not yet settle. A second open question surfaces from the gap between the closed-form linear-quadratic planner and the branched-rollout REINFORCE planner on the cobweb panel. When the environment admits a closed-form planner, branched rollouts pay a variance cost that closed-form planning does not. Whether this gap survives on environments without analytical planners (precisely where deep MBPO is deployed in practice) is the empirical question the deep continuous-control benchmarks were designed to answer. The analogy to E-stability is suggestive rather than formal. The operators, fixed points, and misspecification objects differ, and a stability theory for model-based reinforcement learning in economic environments remains to be built. \citet{Asadi2018Lipschitz} is the closest current analogue to a stability condition for the reinforcement-learning side, stating that value error grows geometrically along the rollout at rate $K_F$ (the Lipschitz constant of the learned dynamics) and remains finite when $\gamma K_F < 1$. But the result is a per-rollout contraction bound on value error rather than a fixed-point characterization of the joint learning dynamics in the sense of \citet{MarcetSargent1989}.

Read against the monograph as a whole, this chapter is where the simulator dependence of reinforcement learning becomes explicit. The question of when a learned model is accurate enough for an economic policy counterfactual is the open one the monograph carries into its conclusion.
